\subsection{Cinemática Inversa}
\label{sec:cinematica_inversa}

A cinemática inversa consiste em determinar o vetor articular que produz uma pose desejada do efetuador em relação à base.
O problema é não linear, pois envolve equações trigonométricas resultantes da composição de transformações homogêneas \cite{craig2005,spong2006robot}.
Como um manipulador 3R possui menos de seis graus de liberdade, não é capaz de realizar combinações arbitrárias de posição e orientação em três dimensões, estando limitado ao subespaço de poses compatíveis com sua geometria \cite{craig2005}.

A existência de solução depende do espaço de trabalho: a pose desejada deve ser alcançável e compatível com a configuração $Z$--$X$--$X$ do \gls{mr}.
Além disso, manipuladores revolutos podem admitir múltiplas soluções articulares para uma mesma pose, exigindo que a escolha final respeite continuidade da trajetória e limites mecânicos \cite{craig2005}.

Para manipuladores com poucos graus de liberdade, não há, em geral, expressão analítica fechada que relacione diretamente a pose desejada aos ângulos das juntas \cite{sciavicco_siciliano2010}.
Neste trabalho adota-se a formulação diferencial, que utiliza a Jacobiana geométrica para relacionar velocidades cartesianas (lineares e angulares) e articulares.
A solução aproximada é obtida pela pseudo-inversa amortecida da Jacobiana, que fornece as velocidades das juntas necessárias para reduzir o erro operacional.
A integração dessas velocidades gera a trajetória correspondente no espaço de juntas, restringindo a solução às poses fisicamente admissíveis pelo manipulador 3R.

\subsubsection{Erro operacional}

A formulação diferencial da cinemática inversa baseia-se na definição do erro no espaço operacional. Seja ${}^{0}\!\vec{x}_{d}$ a posição desejada do efetuador e ${}^{0}\!\vec{x}$ a posição atual obtida pela cinemática direta. Define-se o erro cartesiano como \cite{sciavicco_siciliano2010}:

\begin{equation}
\label{eq:erro_operacional}
\vec{e}
=
{}^{0}\!\vec{x}_{d}
-
{}^{0}\!\vec{x}.
\end{equation}

Esse erro representa a diferença instantânea entre a posição desejada e a posição atual da ferramenta e constitui a variável regulada pela cinemática inversa diferencial.

\subsubsection{Erro de orientação}

Além da posição, é necessário considerar o alinhamento da ferramenta com uma direção desejada no espaço, particularmente na fase de aproximação para acoplamento.
Seja ${}^{0}\!\hat{z}_{E}$ o eixo da ferramenta expresso no referencial da base e ${}^{0}\!\hat{d}$ um versor que representa a direção desejada.
Define-se o erro de orientação como

\begin{equation}
\label{eq:erro_orientacao}
\vec{e}_{\omega}
=
{}^{0}\!\hat{z}_{E}
\times
{}^{0}\!\hat{d},
\end{equation}

de modo que $\vec{e}_{\omega} = \vec{0}$ indica alinhamento perfeito entre o eixo da ferramenta e a direção desejada. Esse erro vetorial é utilizado na construção da componente angular da velocidade operacional desejada.


\subsubsection{Velocidade cartesiana desejada}

A partir dos erros operacional e de orientação, definidos em \eqref{eq:erro_operacional} e \eqref{eq:erro_orientacao}, determina-se a velocidade operacional desejada para o efetuador \cite{sciavicco_siciliano2010}.
Seguindo a formulação clássica de controle cinemático apresentada em \cite{sciavicco_siciliano2010}, utiliza-se uma lei proporcional no espaço operacional, separando as componentes linear e angular

\begin{equation}
\label{eq:velocidade_cartesiana_desejada}
\vec{v}
=
\dot{\vec{x}}_{d}
+
K_{p}\,\vec{e},
\qquad
\vec{\omega}_{d}
=
K_{\omega}\,\vec{e}_{\omega},
\end{equation}

em que $\dot{\vec{x}}_{d}$ é a velocidade cartesiana de referência (que pode ser nula em tarefas de \emph{set-point}), $K_{p}>0$ e $K_{\omega}>0$ são ganhos proporcionais escalares.
A velocidade operacional desejada é reunida no vetor

\begin{equation}
\label{eq:velocidade_operacional_desejada}
\vec{y}_{d}
=
\begin{bmatrix}
\vec{v} \\
\vec{\omega}_{d}
\end{bmatrix}
\in \mathbb{R}^{6},
\end{equation}

que combina o objetivo de posicionamento com o objetivo de alinhamento da ferramenta.


\subsubsection{Lei de cinemática inversa diferencial}

A relação entre a velocidade operacional desejada e as velocidades articulares é obtida por meio da pseudo-inversa amortecida da Jacobiana geométrica. Definindo
\begin{equation}
\mathbf{J}_{g}(\vec q)
=
\begin{bmatrix}
\mathbf{J}_{v}(\vec q) \\
\mathbf{J}_{\omega}(\vec q)
\end{bmatrix}
\in \mathbb{R}^{6\times 3},
\end{equation}

a lei de cinemática inversa diferencial é escrita como

\begin{equation}
\label{eq:ik_diferencial}
\dot{\vec{q}}
=
\mathbf{J}_{\lambda}^{\dagger}(\vec q)\,\vec{y}_{d},
\end{equation}

em que $\mathbf{J}_{\lambda}^{\dagger}(\vec q)$ denota a pseudo-inversa amortecida da Jacobiana geométrica e $\vec{y}_{d}$ é a velocidade operacional desejada definida em \eqref{eq:velocidade_operacional_desejada}.


\subsubsection{Pseudo-inversa amortecida}

A pseudo-inversa amortecida empregada na lei de cinemática inversa diferencial \eqref{eq:ik_diferencial} é definida segundo \cite{sciavicco_siciliano2010} 

\begin{equation}
\label{eq:pseudoinversa_amortecida}
\mathbf{J}_{\lambda}^{\dagger}
=
\mathbf{J}^{T}
\bigl(
\mathbf{J}\,\mathbf{J}^{T}
+
\lambda^{2}\mathbf{I}_{m}
\bigr)^{-1},
\end{equation}

onde $\lambda > 0$ é o fator de amortecimento que regulariza o operador, $\mathbf{I}_{m}$ é a matriz identidade de ordem $m$ (número de coordenadas da tarefa) e $\mathbf{J}$ representa a Jacobiana da tarefa considerada (no caso geral, a Jacobiana geométrica $\mathbf{J}_{g}$). Esse operador garante estabilidade numérica e melhor comportamento na vizinhança de singularidades.


\subsubsection{Integração explícita da lei diferencial}

Uma vez determinada a velocidade articular instantânea $\dot{\vec q}(t)$ pela lei de cinemática inversa diferencial, a trajetória no espaço de juntas é obtida por integração explícita no tempo. Utilizando o método de Euler explícito com passo $\Delta t$, tem-se

\begin{equation}
\label{eq:integracao_explicita_q}
\vec q(t + \Delta t)
=
\vec q(t)
+
\dot{\vec q}(t)\,\Delta t.
\end{equation}

Essa atualização é aplicada iterativamente, de modo que, a cada passo de integração, o erro operacional é recalculado e a lei diferencial de cinemática inversa é novamente avaliada, conforme proposto em \cite{sciavicco_siciliano2010}.

\subsubsection{Condições tratadas na implementação}

A formulação de cinemática inversa diferencial empregada neste trabalho incorpora apenas os mecanismos presentes explicitamente no código.
As seguintes condições são tratadas:

(i) Regularização em regiões próximas de singularidade.  
A pseudo-inversa amortecida é utilizada na forma \cite{sciavicco_siciliano2010}:

\begin{equation}
\mathbf{J}_{\lambda}^{\dagger}
=
\mathbf{J}^{T}
\bigl(
\mathbf{J}\mathbf{J}^{T}
+
\lambda^{2}\mathbf{I}_{m}
\bigr)^{-1},
\end{equation}

de modo que o termo $\lambda^{2}\mathbf{I}$ evita crescimento excessivo das velocidades articulares nas vizinhanças de perda de posto da Jacobiana.

(ii) Correção proporcional no espaço operacional.  
A velocidade desejada no espaço operacional é construída como

\begin{equation}
\vec{y}_{d}
=
\begin{bmatrix}
\dot{\vec{x}}_{d} \\
\vec{0}
\end{bmatrix}
+
\begin{bmatrix}
K_{p}\,\vec{e} \\
K_{\omega}\,\vec{e}_{\omega}
\end{bmatrix},
\end{equation}

de modo a aplicar ação proporcional sobre os erros de posição e de orientação, contribuindo para a estabilidade local da solução diferencial.

(iii) Limitação da velocidade articular.  
A velocidade resultante é saturada a cada componente

\begin{equation}
\dot{q}_{i}
\in
[-\dot{q}_{\max},\,\dot{q}_{\max}],
\end{equation}

de modo a impedir comandos articulares superiores à capacidade do atuador. Não há verificação de limites de posição.

(iv) Estabilidade local da solução.  
A combinação de $K_{p}>0$, $K_{\omega}>0$ e $\lambda>0$ assegura estabilidade local da lei diferencial

\begin{equation}
\dot{\vec{q}}
=
\mathbf{J}_{\lambda}^{\dagger}\,\vec{y}_{d},
\end{equation}

desde que a Jacobiana geométrica mantenha posto pleno na vizinhança da configuração considerada.